
%\needspace{4\baselineskip}
\color{uniblue}{\subsection{Токовая отсечка (ТО)}}
\color{black}

\begin{enumerate}[label=\arabic{section}.\arabic{subsection}.\arabic*, labelsep=4pt, leftmargin=0pt, itemindent=57pt]

\item
Токовая отсечка предназначена для максимально  быстрой  ликвидации  коротких  замыканий, возникающих в начале рабочей зоны.

\item
В устройстве токовая отсечка представлена функцией «ТО». Функция контролирует повышение любого фазного тока ($I_{A\ \textsl{мгн}}, I_{B\ \textsl{мгн}}, I_{C\ \textsl{мгн}}$) выше заданной уставки, при превышении заданной уставки по току через выдержку времени таймера T1 (уставка «Тср») выдается сигнал на отключение трансформатора. При активном состоянии входного сигнала «ТО на сигн» действие токовой отсечки переводится на сигнал.

\item
Ввод функции в работу на этапе параметрирования устройства осуществляется программным переключателем SGF1 «Ввод\_функции». Информация о введенной функции отображается выходным сигналом «ТО: Ввод».

\item
Функция может быть выведена из работы оперативно путем активации сигнала «ОВ ТО», что характеризуется активным состоянием сигнала «ТО: Оперативный вывод» на выходе функции.

\item
Алгоритм работы функции выполнен в соответствии с рисунком \ref{to:img1}. В таблице \ref{to:tbl1} приведены параметры, необходимые для настройки функции.

\vspace{1mm}
\begin{figure}[H]
\centering
\includegraphics[width=0.85\textwidth,height=0.85\textheight,keepaspectratio]{\fbpath/103. ТО ОЛ/_latex/img/to.pdf}
\captionof{figure}{Функциональная схема логики ТО}
\label{to:img1}
\end{figure}


\FloatBarrier % Все float'ы ДО этой строки будут размещены здесь.
\end{enumerate}
